\documentclass[11pt]{article}

\usepackage[utf8]{inputenc}
\usepackage{amsmath}
\usepackage{amssymb}

\title{Effect Isolation for AI Agent Workflows}
\author{}
\date{}

\begin{document}

\maketitle

\section{Introduction}

Modern AI agents increasingly perform actions with irreversible effects:
sending emails, charging credit cards, placing orders, publishing content,
or interacting with external services. Unlike traditional transactions,
such effects cannot always be rolled back. Consider an agent planning a trip under a fixed budget: flight and hotel
together must cost at most 1500 EUR. To reduce latency, the supervisor
runs two subagents in parallel: one books a flight for 900 EUR, while the
other books a hotel for 700 EUR. Both subagents succeed locally. However,
the resulting itinerary violates the global budget constraint. The agent then attempts to recover. The hotel reservation can be canceled,
but the flight was purchased under a non-refundable fare. Thus every
individual tool call was successful, yet the workflow as a whole produced
an externally visible state that cannot be fully repaired.

At the same time, modern agent systems increasingly employ
parallel execution, speculative planning, and recovery after
failures. These mechanisms improve latency and throughput, but they also
introduce new classes of failures in the external world. Classical transaction processing addressed similar challenges through
transaction isolation levels and anomaly-based reasoning~\cite{berenson1995critique,adya1999}.
We argue that agent workflows require a similar abstraction for
external effects.

This direction is related to advanced transaction models, including
Sagas, compensating transactions, escrow, and ACTA
~\cite{garcia1987sagas,korth1990formal,oneil1986escrow,chrysanthis1991acta}.
These models already studied long-running transactions, semantic
compensation, conditional commutativity, and alternative execution
paths. Our goal is not to rediscover these ideas. Rather, we study what
changes when the elementary operations are black-box agent tools, may be
nondeterministic, and may externalize irreversible effects before the
workflow-level decision is known.

\section{Why Agent Workflows Are Concurrent}

Concurrency in agent systems emerges naturally from several sources. Unlike database transactions, these effects interact with an active
external world and may be irreversible.

\paragraph{Subagent decomposition.}
Supervisor-worker architectures decompose tasks into parallel
activities.

\paragraph{Speculation.}
Tree search, best-of-$k$, and branch exploration evaluate multiple
alternative plans simultaneously.

\paragraph{Latency hiding.}
Long-running approvals or external APIs encourage systems to execute
likely actions early.

\paragraph{Recovery and replay.}
Failures and retries introduce repeated execution of effectful actions.

\section{Two Levels of Agent Operations}

We distinguish two abstraction levels L0 and L1. This distinction is inspired by multilevel transaction management, where
high-level transactional guarantees are derived from the properties of
lower-level operations~\cite{weikum1991principles}. Agent workflows
exhibit a similar separation, but with a different boundary. In classical
systems, lower-level operations are typically under the control of the
database system, and their semantic properties can be derived from the
implementation. In agent systems, L0 operations are often black-box tools
behind external APIs. Their properties---such as idempotence,
invertibility, externalization, and determinism---must therefore be
declared, inferred, or validated empirically. These L0 properties then
constrain which isolation guarantees are achievable for the composed L1
workflow.

\subsection{L0: Tool Operations}

L0 consists of individual effectful operations, e.g.:

\begin{itemize}
    \item \texttt{charge()}
    \item \texttt{send\_email()}
    \item \texttt{book\_flight()}
    \item \texttt{publish\_post()}
\end{itemize}

Each operation possesses properties:

\begin{itemize}
\item idempotence,
\item invertibility,
\item externalization,
\item determinism,
\item commutativity.
\end{itemize}

These properties are independent: \texttt{increment} commutes but is not
idempotent, while \texttt{delete} is idempotent but does not commute. They
therefore form a lattice rather than a single scale (e.g.\ a
reversible/irreversible axis). Commutativity is moreover \emph{relational}
(a property of operation pairs) and often \emph{conditional} on state
(escrow-style~\cite{oneil1986escrow}); determinism, absent from classical
concurrency control, is an additional axis that black-box LLM-backed tools
routinely violate.

\subsection{L1: Workflow Operations}

L1 consists of agent tasks, e.g.:

\begin{itemize}
    \item \texttt{HireCandidate}
    \item \texttt{BookTrip}
    \item \texttt{NegotiateContract}
    \item \texttt{PurchaseEquipment}
\end{itemize}

Workflows turn individual tool calls into a logical unit of work. Tool properties constrain the guarantees achievable by workflows, but
isolation itself is defined at L1.

\section{Effect Histories}

We model execution as histories consisting of:

\begin{itemize}
\item $effect(o)$: externalization of operation $o$ --- the moment its
effect becomes visible outside the runtime and is no longer under full
transactional control,
\item $cmp(o)$: compensation,
\item $dep(o_2 \leftarrow o_1)$: $o_2$ causally depends on the effect of
$o_1$ (its decision was made observing $effect(o_1)$),
\item $commit(w)$,
\item $abort(w)$.
\end{itemize}

Each $effect(o)$ carries an outcome status: \emph{confirmed},
\emph{failed}, or \emph{unknown} (no acknowledgment was received). We
additionally define:

\begin{itemize}
\item $resolved(w)$: a $commit(w)$ or $abort(w)$ event has occurred, i.e.\
the workflow's outcome is decided and will not change.
\item $survives(o)$: $effect(o)$ may have occurred (\emph{confirmed} or
\emph{unknown}) and there is no later successful $cmp(o)$ --- the effect
remains in the external world.
\end{itemize}

Unlike classical read-write histories~\cite{adya1999},
effect histories describe what the workflow irreversibly changed in the
external world. The interesting anomalies live in the window between
$effect$ and $resolved$/$cmp$ --- a window absent from read-write models.

\section{Effect Anomalies}

Each anomaly is a \emph{consequence of a missing property}. We give a
one-line scenario and a formal trace.

\paragraph{A1: Duplicated Effect}
A payment request times out, the agent crashes, and after restart the
same invoice is paid again.
\begin{verbatim}
effect(pay_invoice) = confirmed
[crash; replay]
effect(pay_invoice) = confirmed    same logical op, not idempotent
=> invoice paid twice
\end{verbatim}
Non-determinism compounds this: if the operation is also non-deterministic,
replay yields not a duplicate but a \emph{different} effect, and an
\emph{unknown} outcome (A6) cannot even be checked for equality with the
replayed one. Determinism is thus the property that makes recovery by
replay sound.

\paragraph{A2: Irreversible Regret}
A non-refundable ticket is purchased before the remainder of the trip
is finalized, then must be undone.
\begin{verbatim}
effect(book_flight) = confirmed    non-refundable (irreversible)
abort(BookTrip)                    budget violated -> must undo
cmp(book_flight)    = failed
=> survives(book_flight): cannot retract; money lost
\end{verbatim}

\paragraph{A3: Contaminated Speculation}
A committed branch externalizes an action based on a temporary
reservation made by another branch that is later canceled.
\begin{verbatim}
effect(reserve_room R) = confirmed [b1]    b1 holds room R for 14:00
dep(send_invites <- reserve_room R) [b2]   b2 relies on R being held
effect(send_invites) = confirmed [b2]      "meet in R at 14:00" mailed out
commit(b2)                                 b2 wins and is committed
abort(b1); cmp(reserve_room R)             b1 loses; room R released
=> b2 already announced R, which no longer survives
\end{verbatim}
The contamination is causal, not about reversibility: $b2$ trusted an
effect of an unresolved branch. This is the effect-world analogue of a
\emph{dirty read} --- reading uncommitted state that is later rolled
back --- except the ``read'' is a causal dependency on an external effect.
The harm (invitations already sent) would arise even if those invitations
were themselves reversible.

\paragraph{A4: Premature Externalization}
An offer letter is sent before legal approval completes.
\begin{verbatim}
effect(send_offer) = confirmed     letter externalized
NOT resolved(Hire)                 commit/abort(Hire) not yet reached
[later] abort(Hire)                approval denied
=> effect(send_offer) preceded resolved(Hire)
\end{verbatim}
Premature externalization is what separates E2-seq from E2-spec: E2-seq
permits it (the effect may precede the workflow's resolution), while
E2-spec forbids it (effects that may not survive must be staged or gated).

\paragraph{A5: Phantom Compensation}
A supplier receives an offer, reacts to it, and the offer is later
withdrawn. The external world has already incorporated the temporary
state.
\begin{verbatim}
effect(send_offer) = confirmed [A]   supplier sees offer; window opens
effect(react) = confirmed [X]        supplier acts on it
dep(react <- send_offer)             X's action depends on the offer
cmp(send_offer) [A]                  offer withdrawn; window closes
NOT commute(send_offer, react)       cannot reorder: X already acted
=> supplier acted on a now-withdrawn offer; world diverges
\end{verbatim}
This is the only \emph{relational} anomaly (it turns on commutativity of a
pair, not a single operation), which is why it requires a commutativity
\emph{contract} rather than a per-operation hint, and why it is
irremediable when the competitor is the open world (see E3 below).

\paragraph{A6: Orphaned Compensation}
A payment request times out, its status becomes unknown, and the system
issues a refund without knowing whether the original payment succeeded.
\begin{verbatim}
effect(pay_invoice) = unknown      request timed out; outcome unknown
cmp(pay_invoice)    = refund       refund issued blindly
=> if payment never happened, refund is spurious (acting on unknown)
\end{verbatim}
The anomaly is the \emph{blind action} at the moment of $cmp$, not the
eventual outcome; the remedy is to resolve $unknown$ before compensating.

\section{Effect Isolation Levels}
\label{sec:levels}

We conjecture that these anomalies induce a hierarchy of effect isolation
levels, each forbidding a prefix of the catalog. The strict hierarchy
consists of E0, E1, E2-seq, E2-spec, and E3: E1 forbids A1, E2-seq forbids
A1--A3, E2-spec forbids A1--A4, and E3 forbids A1--A6. E2-seq captures
compensation-based sequential execution: it permits premature
externalization (A4) but still excludes contaminated speculation (A3),
requiring that aborted effects leave no surviving residue under the
workflow's compensation assumptions. E2-spec is stricter, additionally
forbidding A4: effects that may not survive must be staged or gated until
the workflow-level decision is known. The anomaly ordering is chosen so
that each level forbids a prefix; the catalog still groups
agent-created anomalies (A1--A4) before externally-caused ones (A5--A6).

\begin{center}
\begin{tabular}{l l p{3.2cm} p{3.6cm}}
\textbf{Level} & \textbf{Forbids} & \textbf{Informally} & \textbf{Where today} \\
\hline
E0 & --- & effects escape immediately & most agent frameworks \\
E1 & A1 & at-most-once effects & ACRFence~\cite{zheng2026acrfence} \\
E2-seq & A1--A3 & compensation-safe effects & RAC~\cite{perera2026rac} \\
E2-spec & A1--A4 & speculation-safe effects & Atomix~\cite{mohammadi2026atomix}; Cordon~\cite{chen2026cordon} \\
E3 & A1--A6 & serializable external history & controlled gateways / no general system yet \\
\end{tabular}
\end{center}

The systems in the last column should be read as implementations of
particular points in the hierarchy under their stated assumptions, rather
than as formal completeness results. ACRFence targets E1 for the
nondeterministic-replay case: after checkpoint-restore an LLM re-synthesizes
a syntactically different but semantically equivalent irreversible call,
which servers accept as new (a duplicate, our A1); ACRFence logs effects and
enforces replay-or-fork to suppress it. (It validates the attack and
proposes the mitigation rather than proving a guarantee.)
RAC targets E2-seq under compensation assumptions: effects externalize
\emph{immediately} during execution --- the runtime does not gate them,
which is precisely why it is E2-seq rather than E2-spec --- and the runtime
records them in a durable log and executes their compensations on abort,
in dependency order so dependents are undone before parents. This excludes
A2 only when the recorded compensations are correct, safe, complete, and
eventually successful; otherwise the workflow may still exhibit irreversible
residue (indeed, when no compensation is found RAC assumes the tool has no
side effect).
Atomix and Cordon implement the stricter E2-spec regime
under a mediated runtime boundary, but by different mechanisms: Atomix
uses frontier-gated settlement under speculation and contention, whereas
Cordon uses task-level transactional containment and validation before
releasing staged effects. None of these systems provides a general E3
guarantee over an active external world.

The difference between E2-seq and E2-spec is whether the runtime controls
the \emph{moment} of externalization. E2-seq allows an effect to become
visible before the workflow outcome is resolved, provided that abort can
remove its residue and no committed branch depends on effects from an
aborted branch. E2-spec is stricter: effects that may not survive must be
staged or gated until the workflow-level decision is known.

These levels are defined for L1 workflows rather than individual tools.
The hierarchy reflects how external the cause of an anomaly is. A1--A4 are
anomalies the agent creates itself and can in principle prevent through
deduplication, gating, dependency tracking, and branch isolation. A5--A6
arise from interaction with an uncontrolled external world: a reacting
competitor, or an unacknowledged external outcome.

Interactive tasks explain why E2-seq is sometimes the best achievable
regime. Consider an agent negotiating on a user's behalf. It cannot learn
the counterparty's reaction without sending an offer, and sending an offer
is already an effect: the action must precede the observation. Full gating
(E2-spec) is therefore impossible. The best attainable guarantee is
E2-seq: the workflow may externalize effects sequentially, but must avoid
irreversible residue in discarded branches. When the inputs needed to
resolve the workflow can be observed without effect, for example through a
\texttt{quote} operation, the stronger E2-spec guarantee becomes
attainable.

E3 should be understood as a mediated serializability target rather than
as a guarantee available in arbitrary open-world settings. It requires
control over the \emph{order} and atomic visibility of external effects:
the system must mediate the order in which effects become visible, resolve
unknown outcomes before compensation, and make dependent effects
atomically visible. Under these assumptions, phantom compensation and
orphaned compensation can be excluded, yielding a serializable external
history.

Without such mediation, E3 becomes an impossibility boundary. The system
may still provide E2-style guarantees over its own actions, but it cannot
in general prevent external actors from observing and reacting to
intermediate states. In open-world settings, the best one can provide is
reconciliation, audit, or application-specific constraints over a history
that the system did not fully control.

\section{Research Directions}

Several questions naturally emerge from this framework.

\subsection{Achievability}

Given:
\[
E^*(W)
\]
the maximum achievable isolation level of workflow $W$.

We conjecture:
\[
E^*(W) = F(P, S, C)
\]
where:

\begin{itemize}
\item $P$: per-tool properties (commutativity, invertibility, idempotence,
externalization, determinism);
\item $S$: workflow structure --- ordering, parallelism and speculation,
the dependency graph, and whether read-only (\texttt{quote}) phases precede
effects;
\item $C$: task constraints and invariants (e.g.\ the budget bound).
\end{itemize}

A subtlety in $S$: an agent's execution graph is not fixed in advance. The
agent decides nondeterministically what to call next, so the graph is
\emph{generated dynamically}. We therefore read $S$ not as a single graph
but as the \emph{space} of histories the agent may produce under a given
tool set and policy, and $E^*(W)$ as the level guaranteed across all of
them --- or, equivalently, the level a runtime mechanism can enforce as
effects appear, rather than by static analysis of a fixed plan (the
approach taken by systems such as Atomix~\cite{mohammadi2026atomix}). This structural nondeterminism
(\emph{which} operation is invoked) is a distinct axis from the operational
nondeterminism already captured in $P$ (whether a given operation is
deterministic).

The decisive argument differs by level: E1 follows mechanically from
per-tool properties (idempotence), whereas E2-spec and above are dominated by
the tool \emph{set} and \emph{composition}. For instance, the budget
invariant of \texttt{BookTrip} depends on flight and hotel prices. If
prices are revealed only by booking, verifying the invariant requires
externalizing irreversible effects --- a cycle (to check, book; to book
safely, check) --- and E2-spec is unachievable. If a read-only
\texttt{quote()} exposes prices without effect, the invariant becomes
checkable before any irreversible effect, and E2-spec becomes achievable. We
call this condition \emph{observability of the invariant}: E2-spec is
achievable iff the values an invariant depends on can be obtained without
externalization, before irreversible effects must occur.

\subsection{Synthesis: the Inverse Problem}

Achievability asks, given a workflow, what level it reaches. The inverse
is often what a developer actually wants: given a target level $E_{target}$
and task constraints $C$, what is the minimal change that makes it
achievable? The change may be a \emph{new tool} (e.g.\ a read-only
\texttt{quote} exposing an invariant's inputs without effect), a
\emph{strengthened tool property} (an idempotency key to reach E1), or a
\emph{restructuring} (inserting a quote-then-commit phase). Each
corresponds to a move in $P$ or $S$, and the output is effectively a
\emph{specification for a missing tool contract} (see the Tool Contracts
direction below).

Two caveats keep the problem honest. Minimality must be defined against
\emph{realistic} extensions: a read-only mirror of an effectful operation
is realistic, whereas making an irreversible effect reversible often is
not. And some targets are unachievable by \emph{any} addition --- e.g.\
interactive tasks where an invariant's input is observable only through the
effect itself (the E2-seq ceiling). The inverse problem must therefore
distinguish ``a tool is missing'' from ``no tool can help.''

\subsection{Commit Spheres}

Some effects cannot be committed independently. In \texttt{BookTrip}, the
charge depends on both bookings (its amount is their sum), and the
confirmation email asserts all three; committing any one without the
others leaves the world inconsistent --- a payment without a booking, or a
confirmation that lies. Such effects must become final together or not at
all. We call a minimal such set a \emph{commit sphere}: the unit of
atomicity for external effects, by analogy with persistence
spheres~\cite{weikum1991principles}, where a minimal set of pages must
reach disk together. (A pure read such as \texttt{search\_flights} has no
effect and belongs to no sphere.)

Commit spheres are the unit of atomicity for external effects, and they
act at two levels. As \emph{all-or-nothing} fixation they already matter
from E2-seq up: a sphere must commit or abort as a whole, so that crashes and
losing branches leave no partial residue (this is failure atomicity). At
E3 a sphere additionally requires atomic \emph{visibility}: its effects
become visible to competitors only all at once, so that no one observes an
intermediate state --- which is exactly what would create phantom
compensation (A5). Their size measures the cost of the guarantee: the
larger the sphere, the more coordination, commit latency, and exposure of
irreversible effects.

Three features distinguish them from the classical
notion: the execution graph is dynamic, so a sphere is discovered
incrementally as dependent effects appear; its elements are heterogeneous
in reversibility, so an irreversible effect cannot be committed on par with
the rest but must be gated; and the dependency graph may contain cycles
(reflection loops).

\subsection{Potential Theorems}

We expect results of the following form.

\begin{itemize}
\item \textbf{Observability is necessary for E2-spec.} If the inputs of a
workflow invariant are observable only through externalizing operations,
then no execution achieves E2-spec: the agent must either externalize
before resolution or fail the task. (This lifts the ``iff'' of the
achievability condition to an impossibility.)

\item \textbf{External ordering precludes E3.} When the order of the
external history is fixed by an actively reacting environment rather than
mediated by the system, phantom compensation (A5) is unavoidable: E3 is
not achievable, and the ceiling is E2-spec over the agent's own actions.
E3 is attainable only when the system mediates the order and atomic
visibility of effects.

\item \textbf{Unknown outcomes preclude exactly-once.} If an effect's
outcome can be \emph{unknown} and durable acknowledgment is unavailable, no
mechanism simultaneously avoids duplication (A1) and orphaned compensation
(A6) --- at-most-once and at-least-once cannot both hold. This is the
agent-level reflection of the classical exactly-once impossibility.

\item \textbf{Commit-sphere lower bound (conjecture).} The size of a forced
commit sphere lower-bounds the coordination cost --- commit latency,
exposure of irreversible effects, and wasted speculative work --- and we
conjecture a protocol that matches it.
\end{itemize}

\subsection{Tool Contracts}

Tool interfaces such as MCP currently expose limited operational
properties. Richer contracts may allow automatic derivation of
achievable isolation guarantees.

\section{Related Work}

Classical transaction models introduced isolation levels and anomaly
definitions~\cite{berenson1995critique,adya1999}. Advanced transaction
models investigated compensation, sagas, semantic recovery, contingency
execution, and dependencies between long-running transactions
~\cite{garcia1987sagas,korth1990formal,chrysanthis1991acta}. Escrow
transactions introduced conditional commutativity and semantic
concurrency control~\cite{oneil1986escrow}.

Recent systems build runtime mechanisms for agent effects; we place them
on the hierarchy in Section~\ref{sec:levels} (ACRFence at E1, RAC at
E2-seq, Atomix and Cordon at E2-spec), and discuss their mechanisms there.
SagaLLM~\cite{chang2025sagallm} applies Saga-style validation to
multi-agent planning, and Stonebraker et
al.~\cite{stonebraker2025consistency} call for isolation-level reasoning in
data-oriented workflow systems. Two systems deserve closer comparison
here. Cordon~\cite{chen2026cordon} introduces semantic transactions ---
task-level boundaries tracking tool intents, result lineage, reversible
state, staged effects, and authority --- which closely mirror our effect
histories, dependency edges, and stage/gate distinction; but it targets
security and least-privilege containment for a single sequential task
rather than graded isolation under concurrency and speculation. On a
complementary axis, S-Bus~\cite{khan2026sbus} formalizes \emph{read}-side
isolation: its Observable-Read Isolation governs what an agent observes of
shared mutable state, where we govern what it irreversibly externalizes ---
the read-side and effect-side of the same consistency problem.

None of these works, however, defines an anomaly-based \emph{hierarchy} of
isolation levels for external effects, or characterizes the achievability
of such levels from tool properties, workflow structure, and task
invariants. Contaminated speculation, phantom compensation, orphaned
compensation, and the joint treatment of read-side and effect-side
isolation remain outside the scope of existing models.

\bibliographystyle{plain}
\bibliography{../../bibliography/flame-stream}

\end{document}